\subsection{The universal covering question remains separate}

\begin{statement}{A.45}{Open}{The full scope of Task A.5}
For arbitrary open \(H\subset\pi\), arbitrary coefficient prime
\(\ell\), and every original cusp packet preserving the indicated
arithmetic cover, we have not determined both the full fixed subgroup
of \(H^{(\ell)}\) under \(s_b(\Delta_H)\) and all its closed
procyclic subgroups with open arithmetic normalizer.
In particular we neither prove nor refute the universal assertion
that both kinds of subgroup are cuspidal. The proved families below
do not form a cofinal family of all open geometric subgroups.
\end{statement}

\begin{proof}[Cases proved and the remaining comparisons]
The actual full profinite fixed subgroup is known for every such
\(H\): it is \(H\cap I\), by A.9. This gives the exact
criterion for a fixed pro-\(\ell\) element to have a fixed lift
to \(H\). It does not compute fixed elements of the quotient
\(H^{(\ell)}\), since the quotient map need not be surjective
on fixed subgroups.

For the quotient problem, A.10 completes all first covers in
Task A.4, including every prime and packet. A.40 completes every
cover with genus-one compactification, with arbitrary isogeny
degree. A.22 also supplies the coefficient prime whenever the
geometric Galois closure is a group of that prime's power order.
The all-prime theorem A.42 supplies the fixed subgroup for any
compact-type reduction with a rational joining path from the
selected cusp to a second cusp. A.43 and A.44 give two explicit
families of unbounded genus satisfying that criterion. A.47 also
completes the normalizer classification at every prime for the
entire odd dihedral family of A.43. A.48 completes every
normalizer in the two-power dihedral family of A.44, including
the previously displayed genus-seven case at coefficient seventeen.
A.49 supplies complete fixed-group and normalizer calculations
for explicit covers in every proper genus, including all even genera;
its extra cusps lie on a rational core in an actual supersingular
Kummer model.
A.50 and A.51 complete both all-prime assertions in the
degree-twenty-seven Heisenberg cover of genus ten. Its residual
surface has actual quaternion multiplication, which supplies
the required reduction and cusp paths at three.
A.52 computes that three-adic reduction directly and completes both
all-prime assertions for every enlargement
\(C\setminus q^{-1}(E[3m])\), for all positive integers \(m\).
Odd multiples use its explicit wild model at three, and even
multiples use the smooth tame map at two.
A.53 proves that every proper compactification occurring in Task A.5
admits infinitely many torsion enlargements with the full all-prime
fixed-inertia theorem. The comparison to the original affine curve is
an open immersion and the covering map is changed by an elliptic
isogeny; this supplies neither a finite refinement over the old map
nor a lift of its fixed elements. Normalizer recognition, when already
known on the original curve, does pass to these enlargements by A.11.
A.54 gives the actual degree-one-hundred-twenty-five Heisenberg cover
of genus fifty-one. Its eight-dimensional multiplicity factor has no CM
subvariety, proved by two persistent global Frobenius fields. All
normalizers are cuspidal for coefficient primes not congruent to one
modulo five; the full fixed group at five follows from A.22.
A.55 completes every remaining split-prime normalizer, including
eleven, using actual cyclic prime-power covers of the genus-two
three-cusp quotient \(w^5=t(t-1)\). Their Pryms have strictly
interior Newton slopes, and an integral transfer calculation detects
even divisible proper-homology classes. The deck translates of
the actual quotient maps detect all remaining quartic-CM lines.
A.58 now completes the full fixed-subgroup comparison at every prime
for this same Heisenberg cover. Its exact five-adic charts give one
elliptic main component, five genus-six and ten genus-two
Artin--Schreier tails. Their genera exhaust fifty-one; all twenty-five
cusps occur in rational clusters away from those tails. Thus the
rational-path criterion A.42 applies to every cusp and packet.

A.56 also gives a normalizer criterion for every proper genus at least two
with potentially good curve reduction at the coefficient prime and
special p-rank at least two. For covers of the original elliptic curve,
at Gaussian split primes it removes the special-rank condition and applies
uniformly to every degree-\(d\) cover when \(p>d\) and
\(p\equiv1\pmod4\). These are normalizer results only.

A.57 completes the special-rank-one case for arbitrary proper curves,
so potential good curve reduction at the coefficient prime is now
sufficient without any rank condition. Consequently every degree-\(d\)
cover has cuspidal normalizers at every \(p>d\), with no congruence
restriction. The outstanding normalizer cases therefore require the
coefficient prime to be at most the degree and the proper covering curve
to have no potentially good reduction at the places above that prime,
unless one of the other proved criteria applies.


A.59 further completes every coefficient-prime normalizer for
every cover whose geometric Galois closure is finite nilpotent,
including all non-Galois covers with such a closure. Its actual
intermediate cover has prime-to-coefficient Galois closure, while
the upper cover has coefficient-prime-group closure and induces
an open pro-prime embedding. No such factorization is asserted
for arbitrary finite groups, and this normalizer theorem does
not determine the full fixed groups of all nilpotent covers.


A.60 completes both assertions at every coefficient prime for
the next actual Heisenberg cover, of degree343, proper genus148
and49cusps. Its explicit seven-adic model has one supersingular
elliptic main component and seven genus21 Artin--Schreier tails;
all marks occur on a rational tree over the origin, disjoint
from those tails. A.61 extends both conclusions uniformly to
every inert-prime Heisenberg cover
\(z^p=\psi_p(x)\), \(p\equiv3\pmod4\), of degree
\(p^3\), genus \(1+p^2(p-1)/2\), and \(p^2\) cusps.
Its divided differential operator proves the local coefficient
identity for all such primes, and the exact charts give one
supersingular elliptic main component and \(p\) tails of genus
\(p(p-1)/2\). This is an unbounded family, with no cofinality
claim and no fixed-group assertion for every nilpotent cover.

A.62 completes the ordinary-prime half by an actual unramified
translation comparison modulo \(p^2\) and separable-isogeny descent.
It gives \(p\) quartic and \(p(p-3)\) quadratic Artin--Schreier
tails, with all cusps on rational clusters of size \(p\).
Thus both assertions are now proved uniformly for the actual
exponent-\(p\) Heisenberg cover at every odd prime \(p\).
This remains a particular unbounded family of geometric subgroups.

A.63 further proves both assertions for every isogeny enlargement
\(C_p\setminus q_p^{-1}(E[pm])\) when \(p\equiv1\pmod4\),
for every \(m\ge1\). An actual noncentral deck subgroup fixes
each relevant special component and supplies a rational companion
for every cusp, including those whose images meet a positive-genus
tail. The comparison to the smaller marked curve is an open immersion.

A.64 removes the congruence restriction in this enlargement conclusion:
every odd Heisenberg prime and every positive enlargement integer now
satisfy both assertions. More generally, if an enlargement integer has
a prime outside the original geometric closure order, the actual tame
good model at that prime gives rational cusp companions and full
fixed inertia at all coefficient primes, including that residue prime.
This includes residue two. It is still an enlargement of the removed
set on the same compactification, with no fixed-lift or cofinality claim.

A.65 also settles every intermediate of each odd Heisenberg cover,
equivalently every cover whose geometric closure has odd-prime exponent
and nilpotency class at most two. The noncentral order-prime quotients
are treated through their actual stable-model maps, rather than by
transferring fixedness through a surjective pro-prime quotient.

A.66 allows an additional abelian geometric factor of order prime to
the Heisenberg prime, including all nonnormal intermediate subgroups.
The actual elliptic-isogeny base change preserves one connected
radicial main and copies each rational marked joining path.
Products of two nonabelian Sylow factors are still outside this
argument; no general nilpotent fixed-group assertion follows.

A.67 supplies a non-nilpotent family: every tetrahedral
\(A_4\) closure and every intermediate has both assertions at
all coefficient primes. The actual three-adic model has a
rational six-cusp centre and four supersingular elliptic
tails. An explicit residue-field argument treats the wild
order-three quotient and its rational two-cusp centre.
This does not extend automatically to coprime elliptic-isogeny
pullbacks: their three-adic dual graph can have cycles.

A.68 completes the two nonstable genus-seven dihedral covers
\(v^2=w^{16}\pm6w^8+1\), and every actual intermediate
of their degree-sixteen closures, at all coefficient primes.
An actual genus-two quotient, exact Igusa invariants and Liu's
residue-two criterion give potential good Jacobians of
\(2\)-rank zero. Explicit principal divisors kill every cusp
difference by eight; norm maps and the Frattini inertia calculation
carry the rational-path criterion to each intermediate.
A.69 now extends this conclusion uniformly to
\(v^2=w^{2m}\pm6w^m+1\) for every two-power \(m\),
including every actual intermediate. The primitive cyclotomic
new factor has special \(2\)-rank zero or is ordinary. The ordinary
alternative contradicts the actual cyclic component action,
finite-flat lifting of smooth fixed points, and the genus of the
rooted marked tree.
A.70 now classifies every actual two-power dihedral closure of
the fixed elliptic base and proves both assertions for all its
intermediates. The three characteristic rotation quotients are
realized by explicit elliptic two-isogenies with distinct dual
kernels; the order-four abelian case is treated separately.
A.71 permits an additional odd-order abelian factor, with every
actual intermediate included. The connected elliptic-isogeny
pullback has a separately proved compact-type model, and its
rational cusp paths lift as literal copies of the original
marked subtree. No special two-rank assertion for the enlarged
Jacobian is needed.
A.72 settles every actual generalized-quaternion two-group
closure and every intermediate. Its central quotient is one
of the classified dihedral covers. A primitive cyclotomic
anti-invariant factor could only have rank zero or be ordinary;
the ordinary case would force a component orbit twice as large
as its total genus, using the quaternion group's unique
involution and the ordinary fourth-root obstruction. The
actual elliptic abelianization quotient gives the required
two-primary cusp divisors.
A.73 supplies the same complete result for every actual
semidihedral two-group closure and its intermediates. Here
the ordinary alternative would force one orbit of elliptic
components; rooting the reduction tree at the whole group's
invariant centre makes them equidistant. Actual fixed-point
lifting and disjoint unmarked child subtrees then contradict
the total genus.
A.74 extends the odd-abelian preservation theorem to all three
two-group families, including every quotient closure and
every intermediate. The same actual connected isogeny
pullback, separately proved tree condition and literal
rational-path lift apply to quaternion and semidihedral
factors as well.
A.75 completes every actual modular two-group closure and
every intermediate. Its unmarked proper stable-tree argument
rules out an ordinary anti-factor by a generically fixed-point-free
outside involution; no marked-fibre slope assertion is assumed.
A.76 combines these four families into the exact class of
two-groups with a cyclic subgroup of index at most two, including
the small orders and every odd abelian extension through the
actual connected isogeny pullback.
A.77 further completes every geometric closure whose commutator
subgroup has at most two elements, with all intermediates and
the automatic odd abelian factors. Actual faithful quotient
curves supply a surjective sum of Jacobian pullbacks, giving
special rank zero for the two-group source. The order-thirty-two
example has genus nine and sixteen cusps and lies outside the
cyclic-index-two class. Products of two nonabelian Sylow factors,
arbitrary two-groups and the universal question remain outside
these conclusions.

A.78 now proves potential good Jacobian reduction of special
rank zero for every prime-power Galois cover of a good
positive-genus rank-zero curve branched at most at one
horizontal section in mixed characteristic. A.79 combines
this with the actual cusp-divisor relation for cyclic-derived
closures: for \(p=2\) or odd \(p\equiv3\pmod4\),
every quotient of a cyclic-derived \(p\)-group times a
prime-to-\(p\) abelian group satisfies both assertions,
including all intermediates. Its actual connected isogeny
pullback preserves each rational companion path.
The rank theorem alone does not give cusp torsion or a
rational companion path for an arbitrary prime-power group;
that fixed-group frontier and the universal question remain
open.


A.80 settles the actual natural genus-two degree-eight intermediate
of the wreath closure \(C_2\wr C_4\), whose derived subgroup
is noncyclic. The remaining marked configuration has component
degrees \((2,0)\), and a cusp can lie on the positive-genus
Prym tail. An exact partial-distance kernel removes the original
elliptic handles while retaining the Prym handles; its non-CM
\(\mathrm{SL}_2\) action then makes every fixed grade Tate.
This settles that actual intermediate at every coefficient prime,
without claiming fixedness for all subgroups of its genus-seventeen
Galois closure.

A.81 proves the same conclusions for any actual genus-two double
cover of an elliptic curve two-primary isogenous to the base,
branched at two isogeny-kernel points and with non-CM elliptic
Prym. In particular all wild two-power torsion enlargements of
the A.80 curve are included, with degree \(2\cdot4^{r+1}\)
and \(2\cdot4^{r+1}-2\) cusps. Each is treated through its
actual degree-two model and two-group closure; no fixed-element
transfer through the affine open immersion is asserted.








A.46 supplies a general
residue-prime normalizer criterion when the potentially good
Jacobian has special p-rank zero. The remaining
normalizer assertions have their own hypotheses: A.11 reduces
a noncuspidal normalized procyclic subgroup to a nonzero proper
homology line, whereas a fixed-group calculation alone makes
no such classification.

The unresolved geometric cases include reductions with a cycle
in their dual graph, and compact-type reductions in which every
path from the selected cusp to a different cusp meets a
positive-genus component. The first-core proof of A.41 uses
positive distances from that joining path to every positive-genus
component. On the path that distance vanishes; for a graph with
cycles the additional graph homology also changes the argument.
Thus one cannot extend A.42 just by allowing zero coefficients
or by applying it to an unspecified semistable alteration.
A.80 supplies a partial-distance replacement for its actual non-CM
elliptic Prym configuration. No general replacement covering all
these remaining cases is proved here.

For normalizers, one must exclude the remaining actual rational
lines in proper homology, not only list their CM characters.
An isogeny factor of the Jacobian does not supply a morphism to
a punctured elliptic target: the inverse image of its missing
point must consist entirely of cusps. A.11 uses that actual
morphism hypothesis. Neither the Jacobian decomposition nor
the absence of new fixed lifts supplies it automatically.

The former example \(Y_8:v^2=w^{16}-1\) at coefficient
seventeen is now settled by A.48 and is not an open case.
The universal problem still includes covers outside all the
proved families and geometric hypotheses above. No cofinality
argument for those families, or least-index determination of
an undecided cover among all open \(H\), has been supplied.
The next required step is an actual normalizer obstruction for
the remaining proper-homology lines, or a constructed noncuspidal
subgroup with its open arithmetic normalizer. Neither an
abstract character nor a vector in a graded completion is such
a counterexample.
\end{proof}
